\documentclass[11pt]{article}
\usepackage[margin=25mm]{geometry}
\usepackage{fontspec}
\setmainfont{Latin Modern Roman}
\usepackage{amsmath,amssymb,mathtools}
\usepackage{unicode-math}
\setmathfont{Latin Modern Math}
\usepackage{xurl,hyperref}
\hypersetup{colorlinks=true,urlcolor=blue}
\setcounter{secnumdepth}{0}
\setlength{\parindent}{0pt}
\setlength{\parskip}{5pt}
\setlength{\emergencystretch}{3em}
\newcommand{\sref}[2]{\hyperlink{source:#1:#2}{[S#2]}}
\newcommand{\cataloguescope}[1]{\paragraph{Recorded target.}#1}
\begin{document}
\section{Kirby Problem 4.127}\label{statement}

\subsection{Definitions and mathematical statement}

\paragraph{Shared notation.}$\mathbb N=\{1,2,\ldots\}$, and $\mathbb Z,\mathbb Q,\mathbb R,\mathbb C$ denote the integers, rationals, reals and complex numbers. Any other conventions are specified locally.


A connected manifold $M$ is aspherical if its universal cover is contractible. For a smooth self-diffeomorphism $f:M\to M$, its mapping torus is
\[
 M_f=(M\times[0,1])/((x,1)\sim(f(x),0)).
\]
Here $M$ has dimension four. The fiber is either closed, or compact with boundary components that admit flat three-dimensional metrics. In the closed case a hyperbolic structure is a metric of constant sectional curvature $-1$ on $M_f$. In the boundary case the hyperbolic structure is interpreted on the interior of $M_f$, with complete finite-volume cusp ends and the boundary regarded as cusp cross sections.
\paragraph{Statement recorded in the supplied source.}
Give necessary and sufficient conditions on $f$ for $M_f$ to admit the indicated hyperbolic structure, for every aspherical four-manifold $M$ in either of these two classes. The criterion must characterize the possible monodromies, not merely construct isolated fibering hyperbolic five-manifolds. \sref{statement}{2}
\subsection{Short English statement}
Characterize the diffeomorphisms of aspherical four-manifolds whose five-dimensional mapping tori are hyperbolic, including the finite-volume cusp case.
\subsection{Sources}
\begin{description}
\item[\hypertarget{source:statement:1}{S1}] ULAM.AI and the UnsolvedMath Contributors, \emph{UnsolvedMath: A Curated Collection of Open Mathematics Problems}, catalogue entry. CC BY 4.0. \url{https://huggingface.co/datasets/ulamai/UnsolvedMath}.
\item[\hypertarget{source:statement:2}{S2}] R. İnanç Baykur, Robion C. Kirby and Daniel Ruberman (eds.), \emph{K3: A New Problem List in Low-Dimensional Topology}, Mathematical Surveys and Monographs 295, American Mathematical Society (2026), Problem 4.127, author preliminary version. \url{https://math.berkeley.edu/sites/default/files/surv-295-ruberman-watermarked-author-pdf.pdf}.
\end{description}
\label{end:statement}
\end{document}
